\section{The Church at the Sickbed: Penance, Extreme Unction, Viaticum, and Commendation}

A Catholic does not cross the threshold of death alone. If Baptism opens the Christian life at the church door and the Eucharist nourishes it through the years, Christ has also provided special sacraments and liturgical rites for the sickroom itself. As the body fails and the adversary redoubles his assaults, the Church comes to the bedside in her priest, bearing absolution, consecrated oil, the living Bread of angels, and the ancient prayers of commendation.

\subsection{The Sacrament of Penance and the Cleansing of the Conscience}

The Church's first care in grave illness is the healing of conscience through Penance. The \emph{Roman Catechism} directs pastors to ensure that the sick person confesses before receiving Viaticum or Extreme Unction, ``as every obstacle which may impede its efficacy should be removed with the greatest care, and as nothing is more opposed to it than a state of mortal guilt.''\endnote{\emph{Catechism of the Council of Trent}, pt.\ 2, On Extreme Unction, trans.\ J.\ Donovan (Baltimore: Lucas Brothers, 1829), p.\ 209; cf.\ \emph{Manual of Prayers for the Use of the Catholic Laity} (New York: Catholic Publication Society, 1889), pp.\ 475, 477.} Where mortal sin has severed divine friendship, sacramental absolution restores sanctifying grace and lifts the dread of unconfessed guilt; even a soul already in grace profits deeply from a humble last confession, entering its final combat stripped of self-deception and clothed in contrition.

\subsection{Extreme Unction: The Anointing of the Sick}

To Penance, Christ added a sacrament instituted specifically for grave bodily affliction: Extreme Unction, or the Anointing of the Sick, promulgated by St.\ James:

\begin{quote}
Is any man sick among you? Let him bring in the priests of the church and let them pray over him, anointing him with oil in the name of the Lord. And the prayer of faith shall save the sick man. And the Lord shall raise him up: and if he be in sins, they shall be forgiven him. (James 5:14--15)
\end{quote}

In its Fourteenth Session (25 November 1551), the Council of Trent defined that Extreme Unction is a true and proper sacrament of the New Law instituted by Christ, intimated in Mark 6:13 and promulgated by St.\ James, which brings both penance and the whole course of Christian life to their consummation and guards the believer at life's close against the final assaults of the adversary. Trent further taught that the sacrament confers the grace of the Holy Ghost, cleanses away any remaining sins together with the remnants of sin, comforts and strengthens the soul with confidence in divine mercy, and restores bodily health whenever it is expedient for the soul's salvation.\endnote{Council of Trent, Session XIV (25 November 1551), Doctrine on the Sacrament of Extreme Unction, Prooemium, chs.\ 1--3, and cans.\ 1--4 (DS 1694--1700, 1716--1719).}

Opening its exposition with Ecclesiasticus 7:40 (``In all thy works remember thy last end, and thou shalt never sin''), the \emph{Roman Catechism} explains why this rite is called the ``last'' unction---the final anointing among those Christ gave His Church, affording ``an easier access to heaven.''\endnote{\emph{Catechism of the Council of Trent}, pt.\ 2, On Extreme Unction, trans.\ Donovan, pp.\ 206--211.} Under the classical pre-1972 \emph{Rituale Romanum} discipline (retained in the 1962 \emph{Rituale Romanum} and printed in the 1889 \emph{Manual of Prayers}), its matter was olive oil consecrated by the bishop, signifying the soothing of anguish, restoration of strength, and light, and its form was spoken by way of prayer over each of the five senses---``Through this holy unction and of His most tender mercy, may the Lord pardon thee whatsoever sins thou hast committed by sight'' (and likewise through hearing, smell, taste and speech, touch, and footsteps)---because bodily recovery rests in divine Providence.\endnote{\emph{Manual of Prayers} (1889), The Sacrament of Extreme Unction, pp.\ 487--488 (\emph{Per istam sanctam unctionem et suam piissimam misericordiam, indulgeat tibi Dominus quidquid per visum deliquisti}), preserving the classical pre-1972 \emph{Rituale Romanum} rite retained in the 1962 \emph{Rituale Romanum}.}

The Catechism warns sharply that ``it is, however, a very grievous sin to defer the Holy Unction until, all hope of recovery now lost, life begins to ebb, and the sick person is fast verging into a state of insensibility'' (p.\ 208). Received while reason still holds sway, its spiritual fruits are profound:

\begin{quote}
It quiets our fear, illumines the gloom in which the soul is enveloped, fills it with pious and holy joy, and enables us to wait with cheerfulness the coming of the Lord, prepared to yield up all that we have received from his bounty, whenever he is pleased to summon us from this world of wo. Another, and the most important advantage derived from Extreme Unction, is, that it fortifies us against the violent assaults of Satan\ldots{} This Sacrament, therefore, arms and strengthens the faithful against the violence of his assaults, and enables them to fight resolutely and successfully against him. (pp.\ 210--211)
\end{quote}

In that classical Roman Ritual discipline, tracing the Cross in blessed olive oil over eyes, ears, nostrils, lips, hands, and feet sealed every sense through which sin had entered. By his Apostolic Constitution \emph{Sacram unctionem infirmorum} (30 November 1972), Pope St.\ Paul VI revised the postconciliar Latin rite so that the sacrament may be conferred with olive oil or, according to circumstances, another suitable plant oil, appoints anointings upon the forehead and hands, and prescribes a single deprecatory essential form asking the Lord, through this holy anointing and His loving mercy, to help the sick person with the grace of the Holy Spirit, free him from sin, save him, and raise him up.\endnote{Pope St.\ Paul VI, Apostolic Constitution \emph{Sacram unctionem infirmorum} (30 November 1972), revising the matter, anointings, and essential form of the postconciliar Latin rite of Anointing of the Sick.} Under both disciplines, the same sacrament heals lingering spiritual languor and prepares the soul to hear ``Blessed are the dead who die in the Lord'' (Apoc.\ 14:13).

\subsection{Holy Viaticum: Food for the Journey}

If Extreme Unction anoints the Christian wrestler for his final contest, the Holy Eucharist received as \emph{Viaticum}---provision for a journey---is the crowning sacrament of the passing itself. Christ promised: ``He that eateth my flesh and drinketh my blood hath everlasting life: and I will raise him up in the last day\ldots{} He that eateth this bread shall live for ever'' (John 6:55, 59; cf.\ 6:54). As the \emph{Roman Catechism} explains, ``Sacred writers also frequently call it `The Viaticum,' as well because it is the spiritual food by which we are supported during our mortal pilgrimage: as also, because it prepares for us a passage to eternal happiness and everlasting glory. Hence, in accordance with the ancient practice of the Church, none of the faithful are suffered to depart this life without being previously fortified with this living bread from heaven.''\endnote{\emph{Catechism of the Council of Trent}, pt.\ 2, On the Sacrament of the Eucharist, trans.\ Donovan, p.\ 147.} So urgent is this precept in danger of death that the Church dispenses the dying from the Eucharistic fast and, in imminent peril, shortens every preparatory prayer so Christ may reach the communicant without delay.\endnote{\emph{Manual of Prayers} (1889), The Holy Viaticum, pp.\ 475--480.}

In the traditional \emph{Rituale Romanum} rite preserved in the 1889 \emph{Manual of Prayers} and the 1962 \emph{Rituale Romanum}, the priest enters the sickroom with the Blessed Sacrament, greets the house---\emph{Pax huic domui} (``Peace be unto this house'')---sprinkles the bed with holy water, hears confession if needed, elevates the Host with \emph{Ecce Agnus Dei}, and joins the sick person in \emph{Domine, non sum dignus}. Then he speaks the classical formula reserved for the threshold of eternity:

\begin{quote}
\emph{Accipe, frater, Viaticum corporis Domini nostri Jesu Christi, qui te custodiat ab hoste maligno, et perducat in vitam aeternam. Amen.} (p.\ 478)
\end{quote}

``Receive, brother, the Viaticum of the Body of our Lord Jesus Christ; may He preserve thee from the wicked enemy, and bring thee unto life everlasting. Amen.'' Afterward he prays that the sacred Body may prove ``an eternal remedy both of body and soul,'' and the devotional thanksgiving places Simeon's joy upon the communicant's lips: ``I hold Thee now, my Love and Sweetness, and will not let Thee go: I gladly bid farewell to the world and all therein; and now I come with joy, my God, to Thee'' (pp.\ 479--480). Carrying the Risen Lord inside his breast, the dying Christian does not walk toward a distant Judge alone.

\subsection{The Commendation of a Departing Soul}

When the last agony arrives, the classical \emph{Rituale Romanum} liturgy of the \emph{Ordo Commendationis Animae}---preserved in the 1889 \emph{Manual of Prayers} and the 1962 \emph{Rituale Romanum}, and retained in its great imperatives \emph{Proficiscere} and \emph{Subvenite} within the simplified postconciliar \emph{Pastoral Care of the Sick}---surrounds the bed with the whole communion of heaven and earth.\endnote{\emph{Manual of Prayers} (1889), Recommendation of a Departing Soul and The Last Agony, pp.\ 511--533, reflecting the traditional \emph{Rituale Romanum} sequence preserved in the 1962 typical edition.} In that traditional order, sprinkling the chamber with holy water, offering the Crucifix to be kissed, and lighting a blessed candle, the priest kneels with the household to pray the Litany for the Dying, pleading through Christ's Birth, Cross, Death, Burial, Resurrection, and Ascension: \emph{A periculo mortis, a mala morte, libera eum, Domine} (``From the peril of death, from an evil death, O Lord, deliver him''). Then, as the soul enters its final passage, the Church speaks her majestic imperative, the \emph{Proficiscere}:

\begin{quote}
\emph{Proficiscere, anima christiana, de hoc mundo, in nomine Dei Patris omnipotentis, qui te creavit; in nomine Jesu Christi, Filii Dei vivi, qui pro te passus est; in nomine Spiritus Sancti, qui in te effusus est\ldots{} hodie sit in pace locus tuus, et habitatio tua in sancta Sion.} (pp.\ 514--515)
\end{quote}

``Go forth, O Christian soul, from this world, in the name of God the Father Almighty, who created thee; in the name of Jesus Christ, Son of the living God, who suffered for thee; in the name of the Holy Ghost, who was poured forth upon thee\ldots{} may thy place be this day in peace, and thine abode in holy Sion.'' Summoning every rank of angels, patriarchs, apostles, martyrs, confessors, and virgins, the prayers that follow (\emph{Deus misericors} and \emph{Commendo te}) beg God to renew all that human frailty corrupted, command Satan and his legions to flee like smoke before the angelic escort, and pray that the gentle and joyful countenance of Jesus Christ appear to the soul (pp.\ 515--518). Next rises the litany of biblical deliverances---\emph{Libera, Domine}, invoking Henoch, Noe, Abraham, Job, Isaac, Lot, Moses, Daniel, the Three Children, Susanna, David, Peter, Paul, and St.\ Thecla---followed by \emph{Commendamus tibi}, begging Christ to recognize His own handiwork formed by ``Thee, the only living and true God'' (pp.\ 519--522).

At the moment of departure, the Church chants the responsory \emph{Subvenite}, summoning the saints and angels to carry the soul before the Most High---``May Christ receive thee, who hath called thee, and may the Angels bear thee into Abraham's bosom''---and closes in the traditional Roman Ritual with \emph{Tibi, Domine, commendamus}: ``To Thee, O Lord, do we commend the soul of Thy servant, that being dead to the world he may live unto Thee; and whatsoever sins he has committed through the frailty of his mortal nature, do Thou, by the pardon of Thy most merciful love, wash away'' (pp.\ 531--533).
